\section{Simultaneous domination and a central defect filler}
\label{cen:section}

The slots of Proposition~\ref{prop:slots} provide a central matrix
cone $\alpha$ with weight $h$, a positive contraction $P=s_1$,
and orthogonal equivalent slots $s_j$.  Put $E=A_\omega$ and
\[
 \Phi(a)=(1-h)a,\qquad \Psi(a)=P^2a\quad(a\in A).
\]
We will extend the comparison relations to obtain
$\Phi(a)\in\Ideal_E(\Psi(a))$ for every $a\in A_+$.
The individual comparison witnesses depend on $a$.  To fill the
defect centrally, for each prescribed finite $F\subset A$ and
$\delta>0$ we need one finite list $v_1,\ldots,v_N\in E$ such that
\[
 \left\|\Phi(a)-\sum_{j=1}^N v_j^*\Psi(a)v_j\right\|<\delta
 \qquad(a\in F).
\]
The first part of this section proves this implication for weighted
inclusions.  Nuclearity allows us to pass from pointwise ideal
containment to an inclusion of support-representation kernels; the
two tensor-ideal lemmas establish the needed bridge.

The slots then pack the common sum into $z=\sum_j s_jv_j$.
These elements need not be uniformly bounded or central.
Approximate-unit tests will give bounded compressions with
commutators tending to zero, and hence a bounded central element
$S$ satisfying
\[
 \alpha(e_{11})S=S,\qquad S^*S a=(1-h)a\quad(a\in A).
\]
This defect-filling argument is related to Matui--Sato's use of
excision in small central sequences and property \textup{(SI)}
\cite{MatuiSato2012}.

\subsection{Tensor ideals and support representations}

\begin{lemma}\label{cen:elementary-tensor}
Let $D$ and $C$ be $C^*$-algebras.  Every nonzero closed two-sided ideal
of $D\otimes C$ contains a nonzero tensor $a\otimes b$ with
$a\in D_+$ and $b\in C_+$.
\end{lemma}

\begin{proof}
Let $0\ne x\in J_+$.  There is a pure state $\lambda$ of the
unitization $D^{\sim}$ such that
$b=(\lambda\otimes\id_C)(x)\ne0$.
Indeed, product positive functionals detect nonzero positive elements of
a spatial tensor product.  Fixing the second functional gives a nonzero
positive affine function on the state space of $D^{\sim}$, which is
nonzero at some extreme point.  The resulting pure state does not
annihilate $D$.  Its GNS representation is therefore nondegenerate on
$D$; write $\xi$ for its unit cyclic vector.

Lemma~\ref{norm:excision} gives,
on any prescribed finite set of coefficients in $D^{\sim}$, positive
contractions $g\in D^{\sim}$ with
\[
 \lambda(g)=1,\qquad
 \|gcg-\lambda(c)g^2\|\longrightarrow0.
\]
Here $\pi_\lambda(g)\xi=\xi$, since $g$ is a positive contraction.
Choose once and for all a positive contraction $u\in D$ such that
$\|\pi_\lambda(u)\xi-\xi\|<1/8$.  Approximation of $x$ by finite
algebraic sums and excision of their first coefficients show that
\[
 \dist\bigl((ug^2u)\otimes b,J\bigr)\longrightarrow0.
\]
The element used to approximate it is
$(ug\otimes1)x(gu\otimes1)$, which belongs to $J$; all products with
$1$ here are multiplier products.  Moreover,
\[
 \|ug^2u\|\ \ge\
 \|\pi_\lambda(g)\pi_\lambda(u)\xi\|^2
 \ > (7/8)^2>1/2.
\]
Thus the first factor retains a fixed norm lower bound as the excision
accuracy improves.

We use the following elementary cut observation.  If $a,b$ are positive
and $\|a\otimes b+J\|<t^2$, then
\begin{equation}\label{cen:product-cut}
 (a-t)_+\otimes(b-t)_+\in J.
\end{equation}
To see this, represent $(D\otimes C)/J$ faithfully and nondegenerately.
The commuting multiplier operators representing $a\otimes1$ and
$1\otimes b$ have product of norm less than $t^2$.  Their joint spectrum
has no point with both coordinates larger than $t$.  Joint continuous
functional calculus proves \eqref{cen:product-cut}.
Now fix $0<t<\min\{1/2,\|b\|\}$ and choose $g$ so that the preceding
distance is less than $t^2$.  Both cuts in
$(ug^2u-t)_+\otimes(b-t)_+$ are nonzero, and their tensor product
belongs to $J$.
\end{proof}

\begin{lemma}\label{cen:nuclear-tensor-ideals}
If $D$ is nuclear and $C$ is arbitrary, every closed two-sided ideal
$J\subset D\otimes C$ is generated by the positive elementary tensors
that it contains.
\end{lemma}

\begin{proof}
Let $J_0$ be the ideal generated by those tensors.  If $J_0\ne J$,
primitive ideal separation gives a nondegenerate irreducible
representation $\pi$ with $\pi(J_0)=0$ and $\pi(J)\ne0$.
Let $\pi_D$ and $\pi_C$ be its commuting leg representations, and put
$I=\ker\pi_D$, $L=\ker\pi_C$.  Nuclearity of $D$ and of its quotient
$D/I$ \cite[Corollary~4]{CE1977}, together with the quotient identity for the maximal tensor
product, gives
\begin{equation}\label{cen:tensor-quotient}
 q:D\otimes C\longrightarrow (D/I)\otimes(C/L),\qquad
 \ker q=\overline{I\otimes C+D\otimes L}.
\end{equation}
For completeness, the maximal quotient identity follows by the
universal property of commuting representations: a representation
annihilating the indicated two leg ideals is exactly a pair factoring
through $D/I$ and $C/L$.  The nuclearity of the first leg identifies
both maximal products in this statement with spatial products.

The representation induced by $\pi$ on the target of $q$ is faithful.
Here is the needed verification.  Write
$M=\pi_D(D)''$ and $N=\pi_C(C)''$.  The center of $M$ commutes with
both legs, hence with the irreducible representation $\pi$; it consists
of scalars.  Thus $M$ is a factor, and the same holds for $N$.
If nonzero positive elements $x\in M$, $y\in N$ had $xy=0$, their
support projections $e\in M$, $f\in N\subset M'$ would satisfy
$ef=0$.  Since $f$ commutes with $M$, it annihilates every $vev^*$,
$v\in\mathcal U(M)$.  The join of these projections is the central
support of $e$, which is $1$ in the factor $M$.  This would give $f=0$,
a contradiction.  The leg representations of $D/I$ and $C/L$ are
faithful, so the induced tensor representation kills no nonzero positive
elementary tensor.  Lemma~\ref{cen:elementary-tensor} applied to its
kernel proves faithfulness.

The nonzero ideal $q(J)$ therefore contains a nonzero positive tensor
$\bar a\otimes\bar b$.  Choose positive lifts $a\in D$, $b\in C$.
Since $a\otimes b\in J+\ker q$, compression by
$(1-e)\otimes(1-f)$, where $e$ and $f$ run through positive contractive
approximate units of $I$ and $L$, gives
\[
 \dist(a_{e}\otimes b_f,J)\longrightarrow0,
 \quad a_e=(1-e)a(1-e),\quad b_f=(1-f)b(1-f).
\]
Indeed these compressions send each term of a finite sum from
$I\otimes C+D\otimes L$ to zero in norm, and they preserve $J$.
The quotient images of $a_e,b_f$ remain $\bar a,\bar b$.  Fix
$0<t<\min\{\|\bar a\|,\|\bar b\|\}$.  Once the displayed distance
is less than $t^2$, \eqref{cen:product-cut} yields
$(a_e-t)_+\otimes(b_f-t)_+\in J_0$ with nonzero image under $q$.
Faithfulness of the induced representation contradicts $\pi(J_0)=0$.
\end{proof}

\subsection{The simultaneous approximation}

Gabe's approximate domination theorem
\cite[Definition~3.1 and Theorem~3.3]{Gabe2020} applies to an exact
$C^*$-algebra $D$, an arbitrary $C^*$-algebra $E$, and nuclear
completely positive maps $\Psi,\Phi:D\to E$, with $\Psi$ of order
zero.  It identifies the condition
$\Phi(a)\in\Ideal_E(\Psi(a))$ for every $a\in D_+$ with the
following finite-set approximation: for every finite $F\subset D$
and $\delta>0$, there are $N\in\N$ and $v_1,\ldots,v_N\in E$
such that
\begin{equation}\label{cen:domination-equation}
 \left\|\Phi(a)-\sum_{j=1}^Nv_j^*\Psi(a)v_j\right\|<\delta
\qquad(a\in F).
\end{equation}
Nuclearity is required of the maps; no nuclearity or separability
assumption is imposed on $E$.  We prove the weighted-inclusion case
used here, so that the passage from ideal containment to one common
finite list of compression coefficients is explicit.

\begin{lemma}[Simultaneous domination for weighted inclusions]
\label{cen:weighted-domination}
Let $D\subseteq E$ be a nuclear $C^*$-subalgebra of a $C^*$-algebra
$E$, and let $v,w$ be positive contractions in $\widetilde E\cap D'$.
Put $\Phi(a)=va$ and $\Psi(a)=wa$ for $a\in D$; these products
belong to $E$. Suppose that
\[
 \Phi(a)\in\Ideal_E(\Psi(a))\qquad(a\in D_+).
\]
Then for every finite set $F\subset D$ and every $\delta>0$, there
are an integer $N\ge1$ and $v_1,\ldots,v_N\in E$ such that
\[
 \left\|va-\sum_{j=1}^N v_j^*wa v_j\right\|<\delta
 \qquad(a\in F).
\]
\end{lemma}

\begin{proof}
Let $\mathcal Q$ be the cone of maps
$a\mapsto\sum_{j=1}^N b_j^*wa b_j$, where $b_j\in E$.
It is convex: sums concatenate the lists, and multiplication by a
nonnegative scalar is absorbed into the $b_j$.  It suffices first to
place $\Phi$ in its point-weak closure.

Fix $a_1,\ldots,a_m\in D$ and $\ell_1,\ldots,\ell_m\in E^*$.
We may represent all these functionals in one nondegenerate cyclic
representation $\sigma:E\to\mathcal B(H)$, with cyclic vector $\xi$,
in the form
\begin{equation}\label{cen:common-functional-model}
 \ell_i(b)=\langle\sigma(b)c_i\xi,\xi\rangle,
 \qquad c_i\in\sigma(E)'.
\end{equation}
To verify this, decompose the real and imaginary parts of each
$\ell_i$ into differences of positive functionals.  A positive multiple
of their sum dominates all the positive summands.  In its GNS
representation a dominated positive functional gives a bounded positive
form on vectors $\sigma(b)\xi$; its representing operator commutes
with $\sigma(E)$, as is seen by moving $\sigma(x)$ between the two
arguments of the form.  Linear combinations give the $c_i$ in
\eqref{cen:common-functional-model}.  Extend $\sigma$ unitally to
$E^{\sim}$ when necessary.

Put $C=C^*(1,c_1,\ldots,c_m)\subset\sigma(E)'$.  The commuting
representations define
\[
 \pi:D\otimes C\longrightarrow\mathcal B(H),\qquad
 \pi(a\otimes c)=\sigma(a)c;
\]
continuity for the spatial tensor norm follows from nuclearity of $D$.
For $q=v,w$, let $p_q$ be the support projection of $\sigma(q)$ and
let $\pi_q$ be the restriction to $p_qH$.  This is a reducing subspace
for $\pi$: each $q$ commutes with $D$, and each element of $C$ commutes
with $E^{\sim}$.

We claim
\begin{equation}\label{cen:support-kernel}
 \ker\pi_w\subset\ker\pi_v.
\end{equation}
Suppose $a\otimes c\in\ker\pi_w$.  Commutation with the weight
and positivity give
\[
 \sigma(wa^*a)c^*c=0.
\]
Since $c^*c$ commutes with all of $\sigma(E)$, this identity persists
after multiplying $\sigma(wa^*a)$ on either side by elements of
$\sigma(E)$ and taking closed linear spans.  The ideal hypothesis
$va^*a\in\Ideal_E(wa^*a)$ therefore gives
$\sigma(va^*a)c^*c=0$.  This says
$\pi(a\otimes c)\sigma(v)^{1/2}=0$, and, by density of the range of
$\sigma(v)^{1/2}$ in $p_vH$, gives $a\otimes c\in\ker\pi_v$.
Lemma~\ref{cen:nuclear-tensor-ideals} says that the elementary tensors
in $\ker\pi_w$ generate that kernel, and proves
\eqref{cen:support-kernel}.

Let $K=\ker\pi_w$.  The vectors
\[
 \eta_b=\sigma(w)^{1/2}\sigma(b)\xi,\qquad b\in E,
\]
are dense in $p_wH$: cyclicity first gives density before application
of $\sigma(w)^{1/2}$, and its range is dense in its support space.
Finite sums of their vector functionals have weak-* closure equal to
the positive functional cone of $(D\otimes C)/K$.  Here is the
separation argument.  A selfadjoint element of this quotient that is
nonnegative on all these functionals has a positive representing
operator on $p_wH$, by the density of the $\eta_b$.  Faithfulness of
the quotient representation makes the element positive.  Thus the
polar of the cone of these vector functionals is exactly the positive
cone of the quotient.  The bipolar form of Hahn--Banach separation
identifies its weak-* closure with all positive functionals on that
quotient.  This argument also applies to a nonunital quotient.

The positive functional
\[
 \gamma_v(x)=
 \langle\pi(x)\sigma(v)^{1/2}\xi,
                   \sigma(v)^{1/2}\xi\rangle
\]
factors through this quotient by \eqref{cen:support-kernel}.  Hence,
simultaneously on the finitely many tensors $a_i\otimes c_i$, it is
approximated by finite sums of vector functionals for $\eta_{b_j}$.
On these tensors the identities are
\begin{align*}
 \gamma_v(a_i\otimes c_i)&=\ell_i(va_i),\\
 \sum_j\langle\pi(a_i\otimes c_i)\eta_{b_j},\eta_{b_j}\rangle
   &=\ell_i\left(\sum_jb_j^*wa_i b_j\right).
\end{align*}
They prove the required point-weak closure, with the same finite list
of coefficients for all the tests.  Finally, for a fixed finite list
$a_1,\ldots,a_m$, the set
\[
 \{(T(a_1),\ldots,T(a_m)):T\in\mathcal Q\}\subset E^m
\]
is convex.  Its weak and norm closures coincide by Hahn--Banach
separation.  This proves \eqref{cen:domination-equation}.
\end{proof}

\subsection{From slots to a bounded central element}

\begin{proposition}\label{prop:central-filler}
Let $A$ be a separable nuclear $C^*$-algebra satisfying
\textup{(C)}.
For each $p\ge2$ there are a \cpc order-zero map
$\alpha:M_p\to A_\omega\cap A'$, with $h=\alpha(1)$, and an element
$S\in A_\omega\cap A'$ such that
\begin{equation}\label{cen:filler-relations}
 \alpha(e_{11})S=S,\qquad S^*S a=(1-h)a\quad(a\in A).
\end{equation}
One can choose $\|S\|\le2$.  In the second equation the expression
$(1-h)a$ is computed in the unitization of $A_\omega$ and belongs to
$A_\omega$.
\end{proposition}

\begin{proof}
Fix the map and slots from Proposition~\ref{prop:slots}.  Thus
$P=s_1\ge0$, all $s_j$ are central contractions,
\begin{equation}\label{cen:slot-recall}
 s_i^*s_j=\delta_{ij}P^2,\qquad
 \alpha(e_{11})s_j=s_j,
\end{equation}
and, for every element $a$ of the countable dense family
$\mathcal A$ of positive contractions fixed in
Section~\ref{sec:comparison}, closed under positive rational cuts, and
every rational $\eps>0$, there is a bounded element $r_{a,\eps}$
with
\[
 ((1-h)a-\eps)_+=r_{a,\eps}^*(P^2a)r_{a,\eps}.
\]
Set $E=A_\omega$, $v=1-h$, $w=P^2$, and define
$\Phi(a)=va$, $\Psi(a)=wa$.  Both maps are completely positive
and contractive, and $\Psi$ has order zero.  Since $A$ is nuclear,
composing finite-dimensional completely positive approximations to
$\id_A$ with either map proves that both maps are nuclear.

Letting $\eps\downarrow0$ in the last comparison gives
$\Phi(d)\in\Ideal_E(\Psi(d))$ for $d\in\mathcal A$.  We verify
the extension to an arbitrary positive contraction $a\in A$.  Choose
$d_n\in\mathcal A$ and positive rational $\delta_n\downarrow0$
with $\|d_n-a\|<\delta_n$.  The cuts
$b_n=(d_n-\delta_n)_+$ belong to $\mathcal A\cap\Ideal_A(a)$,
by positive-element perturbation, and $b_n\to a$ in norm.
Centrality of $w$ implies
\begin{equation}\label{cen:ideal-extension}
 w\Ideal_A(a)\subset\Ideal_E(wa):
 \quad w(xay)=x(wa)y\quad(x,y\in A).
\end{equation}
The same assertion follows for the closed ideal by approximation with
finite sums.  Consequently $\Psi(b_n)\in\Ideal_E(\Psi(a))$, and
\[
 \Phi(b_n)\in\Ideal_E(\Psi(b_n))
             \subset\Ideal_E(\Psi(a)).
\]
Taking the norm limit proves the ideal hypothesis for $a$; scaling
proves it for every positive element.

Apply Lemma~\ref{cen:weighted-domination}.  Any resulting sum packs
into a single element
$z=\sum_{j=1}^N s_jv_j\in E$, because, for $a\in A$,
\begin{equation}\label{cen:packed-sum}
 z^*az=\sum_{i,j}v_i^*s_i^*as_jv_j
      =\sum_jv_j^*(P^2a)v_j.
\end{equation}
There is no uniform bound here on $z$ or on the individual $v_j$,
and the coefficients need not commute with $A$.  A compression by
one approximate-unit element will supply both the bound needed for
diagonalization and, through simultaneous tests, control of the
commutators.

Fix a finite set $F$ of contractions in $A$, and let $\eta,\delta>0$.
Choose a positive contraction $t\in A$ from an approximate unit so
that, for all $a\in F$,
\[
 \|t^2a-a\|<\eta,\qquad
 \|tat-a\|<\eta,\qquad
 \|ta^*at-a^*a\|<\eta.
\]
Approximate $\Phi$ to within $\delta$, using one sum in
\eqref{cen:packed-sum}, on
\[
 \{t^2\}\cup\{tat,ta^*at:a\in F\}.
\]
Put $S'=tz$.  The $t^2$ test gives
\begin{equation}\label{cen:bounded-compression}
 \|S'\|^2=\|z^*t^2z\|\le\|vt^2\|+\delta\le1+\delta.
\end{equation}
Also $\alpha(e_{11})S'=S'$ exactly, since $\alpha(e_{11})$
commutes with $t$ and fixes all the $s_j$.  Writing
$\kappa=\delta+\eta$, the three tests give
\begin{align}
 \|S'^*S'a-va\|&<\kappa,\label{cen:first-test}\\
 \|S'^*aS'-va\|&<\kappa,\label{cen:second-test}\\
 \|S'^*a^*aS'-va^*a\|&<\kappa.\label{cen:third-test}
\end{align}
For the first line, multiply
$\|S'^*S'-vt^2\|<\delta$ by $a$; the other two lines follow
directly from their respective tests.  Taking the adjoint of
\eqref{cen:second-test} controls $S'^*a^*S'$ as well.
Now expand the full commutator square:
\begin{align*}
 (aS'-S'a)^*(aS'-S'a)
   ={}&S'^*a^*aS'-S'^*a^*S'a\\
     &-a^*S'^*aS'+a^*S'^*S'a.
\end{align*}
Its four terms are within $\kappa$ of, respectively,
$va^*a$, $-va^*a$, $-va^*a$, and $va^*a$, because $v$ commutes
with $A$.  Therefore
\begin{equation}\label{cen:commutator-bound}
 \|[a,S']\|\le2\sqrt\kappa\qquad(a\in F).
\end{equation}

Take increasing finite sets dense in the unit ball of $A$ and choose
$\eta,\delta\downarrow0$, always with $\delta<1$.  Each resulting
$S'$ has norm at most $\sqrt2$ by
\eqref{cen:bounded-compression}, and may be represented by a sequence
bounded by $2$.  Fix representatives of the finitely many matrix
entries of $\alpha$ throughout.  The countable diagonalization of
Lemma~\ref{lem:diagonal}, applied to these fixed representatives and
the norm tests \eqref{cen:first-test},
\eqref{cen:commutator-bound}, and
$\alpha(e_{11})S'=S'$, gives a single element $S$ of norm at most $2$.
It commutes with every element of $A$, retains the fixed support
equation, and satisfies \eqref{cen:filler-relations} on a dense set,
hence on all of $A$ by continuity.  Only the uniformly bounded $S'$
are diagonalized at this last step; the compression coefficients and
their unbounded sums are not retained as variables.
\end{proof}
